N=60 revivals appear for $5$ and $10$ tori, with reflection parameters
\begin{equation}
\rho = \frac{1}{8}\bigl(5-\sqrt{5}\bigr) \approx 0.276393, \qquad
\rho = \frac{1}{8}\bigl(3+\sqrt{5}\bigr) \approx 0.723607 = 1-\rho .
\end{equation}
The phase parameters take values in
\begin{equation}
\delta \in \left\{\,0,\ \frac{2}{5},\ \frac{4}{5}\,\right\},
\end{equation}
suggesting a phase periodicity of $2\pi/5=1.2566$, or equivalently
$\delta\in\{0,\,2\pi/5,\,4\pi/5\}$. The critical observation is that the
$5\times5$ torus alone contains thirty-six parameter sets yielding
period 60. This is the largest revival period in the entire dataset,
and it corresponds to the smallest torus in the family (in the sense
that $L_x=L_y=5$ is prime). This strongly suggests an inverse
relationship between torus size and maximum revival period for a given
dimension, and it identifies the $5\times5$ torus as the richest single
source of long-period revivals in the separable-coin setting.

The largest dimension of torus which revivals occur is
$10 \times 10$, $N=60$ (4 revivals). In this case, $\rho_x, \rho_y$ are
related to the golden ratio:

Given, $\rho_x = \frac{5 + \sqrt{5}}{10}, \delta=0$.
$\varphi = \frac{1+\sqrt{5}}{2}$
We have:
\begin{equation}
    \rho_x = \frac{1}{\varphi + 2}
    = \frac{1}{\varphi\sqrt{5}}
    = \frac{5-\sqrt{5}}{10}
    \approx 0.276393.
    \label{eq:rho-golden}
\end{equation}
and then $\rho_y = 1 - \rho_x$ and vice versa. In this case, revivals
require \(\delta_x, \delta_y = 0\). A nonzero coin phase
\(e^{i\delta}\) makes the coin eigenvalues irrational.

For the \textbf{$8 \times 8$} torus, All 16 revivals occur at $N=24$
with combinations of $\delta_x, \delta_y \in \{0, \pi/2, \pi, 3\pi/2\}$
at $\rho_x = \rho_y = 0.5$ which reflects the 4-fold symmetry of the
Hadamard coin.
